\chapter{Design}
Design is the backbone of a visualisation---it communicates how the chart uses (which) marks and channels to communicate the information.
Good design may be aesthetic, useful, honest, and minimal\footnote{No, I don't want to list \emph{all} of Dieter Rams \emph{Principles of Good Design}}.

Choosing a design can be a tough iterative process, that cycles between data-near analysis (what does this graph require?), and improvement (how can I make this graph look better?).
The goal of this chapter is to formalize a critique of visualisation design.
Edward Tufte defines two main goals for designs.

\begin{enumerate}
    \item Practice: Graphical Integrity
    \item Theory: Graphical Excellence
\end{enumerate}

\section[Graphical Integrity]{Graphical Integrity \\ {\large Visual Representations of Data Must Tell the Truth}}
There are different ways to lie in a graph, one may offset a scale or omit it completely from the graph.
One can also cleverly remove some data from the dataset to change the meaning of the chart.

To see through some of these distortions, Tufte developed the \emph{lie factor}---a metric which compares the size of effect in the graph with the size of effect in the data.
It is calculated in the following way---where \(S_\text{graph}\) is calculated by dividing the difference between minimal and maximal size in the graph by the minimal size.\footnote{It is important, that the \emph{size} in a graph is relative to the way the data is represented---is it length, area, volume, or something entirely else?}
\(S_\text{data}\) is analogously calculated with minimum and maximum of the data mapped to the graph.
\[\text{LF} := \frac{S_\text{graph}}{S_\text{data}}\]

\section{Graphical Excellence}
Tufte invented a method on how to achieve graphical excellence based on certain design principles.
\begin{description}
    \item[Show Data] Above all else, show the data.

    \item[Prioritise Data Ink] Maximize the share of data ink\footnote{Ink used to display data} within reason.
        \emph{Within reason} because removing e.g., axis labels obfuscates the graph.
        The \emph{data ink ratio} is calculated by dividing the amount of data ink by the total amount of ink in the graph.
        Ideally and naively 0.

    \item[Avoid Chartjunk] Remove unnecessary visual elements in charts that distract the viewer from the information.

    \item[Erase Redundant Data Ink] If some data ink is used multiple times for the same information, remove the duplicates.

    \item[Revise and Edit] Again, having an iterative design process is important for good design.

    \item[Maximize Data Density] Try to display a lot of data in a small area.
        The data density of a graphic is calculated by dividing the number of data entries displayed by the area of the graphic.
        A good strategy is to display multiple visually similar small graphs near each other.
\end{description}
